\documentclass[UTF8,fontset=none,aspectratio=169,10pt]{beamer}
\usepackage{ads-beamer-theme}
\title{高级分布式系统}
\subtitle{第十二讲：检查点、回滚与恢复}
\author{陈鹏飞\quad 计算机学院}
\institute{中山大学}
\date{2025--2026 学年}

\begin{document}
\begin{frame}[plain]
\begin{tikzpicture}[remember picture,overlay]
  \fill[ADSBlue] (current page.north west) rectangle ([xshift=2.4cm]current page.south west);
  \fill[ADSOrange] ([xshift=2.4cm]current page.north west) rectangle ([xshift=2.55cm]current page.south west);
  \node[anchor=west,align=left] at ([xshift=3.3cm,yshift=1.0cm]current page.west) {\fontsize{32}{38}\selectfont\color{ADSBlue}\bfseries 高级分布式系统};
  \node[anchor=west,align=left] at ([xshift=3.3cm,yshift=-0.4cm]current page.west) {\fontsize{21}{27}\selectfont\color{ADSDarkBlue}第十二讲\quad 检查点、回滚与恢复};
  \node[anchor=west,align=left] at ([xshift=3.3cm,yshift=-2.2cm]current page.west) {\large\color{ADSGray}故障后回到哪里，才能让所有进程重新“对得上账”？};
  \node[anchor=east] at ([xshift=-0.8cm,yshift=-3.2cm]current page.east) {\small\color{ADSGray}陈鹏飞\quad 中山大学\quad 2025--2026};
  \draw[ADSBlue,line width=0.8pt] ([xshift=3.3cm,yshift=-1.35cm]current page.west) -- ([xshift=12.6cm,yshift=-1.35cm]current page.west);
\end{tikzpicture}
\end{frame}

\sectionframe{一、恢复的核心问题：一致地回到过去}{检查点保存局部状态，但消息建立了跨进程依赖；单独回滚很容易造成新的错误。}

\begin{frame}{本讲学习目标}
\begin{enumerate}
  \item 定义一致全局状态、全局检查点和恢复线；
  \item 理解 orphan、lost、delayed、duplicate 消息；
  \item 比较非协调、协调、通信诱导和日志型恢复；
  \item 解释 domino effect、输出提交和确定性重放；
  \item 把 Checkpoint/Recovery 映射到 AI 训练和 Kubernetes 工作流。
\end{enumerate}
\vfill
\callout{主线}{恢复不是“每个进程找最近备份”，而是寻找一个在某次无故障执行中可能出现的一致切。}
\end{frame}

\begin{frame}{局部检查点与全局检查点}
\begin{itemize}
  \item $C_{i,k}$：进程 $P_i$ 的第 $k$ 个局部检查点；
  \item 每个进程在开始执行前至少有 $C_{i,0}$；
  \item 全局检查点是一组“一进程一个局部检查点”；
  \item 一致全局检查点必须能在无故障执行中出现，不能让一个进程“看到”尚未由发送者发出的消息。
\end{itemize}
\end{frame}

\begin{frame}{一致性条件}
\begin{itemize}
  \item 若接收事件出现在恢复状态中，则对应发送事件也必须出现在恢复状态中；
  \item 不允许出现 orphan message：接收被保留，但发送已被回滚；
  \item 可以存在 in-transit message：发送已发生，接收尚未发生；
  \item 恢复线本质上是一条一致切：所有进程沿自己的本地事件序列回到同一条全局历史。
\end{itemize}
\vfill
\callout{与快照的联系}{快照记录“现在的全局状态”；恢复线选择“过去的全局状态”，两者都依赖一致切。}
\end{frame}

\begin{frame}{消息在恢复中的五种命运}
\begin{tabularx}{\linewidth}{>{\bfseries}p{2.6cm}X}
\toprule
In-transit & 发送已记录，接收尚未记录；恢复后仍可继续到达。 \\
Lost & 发送记录被保留，但接收事件因回滚被撤销；需要重发或日志恢复。 \\
Delayed & 消息到达时间跨过恢复线，接收进程暂时看不到。 \\
Orphan & 接收记录保留，但发送记录被回滚；违反一致性。 \\
Duplicate & 日志重放或重试让同一逻辑消息再次交付；需要序号/幂等。 \\
\bottomrule
\end{tabularx}
\end{frame}

\begin{frame}{一个简单的 orphan 反例}
\begin{enumerate}
  \item P1 在检查点 $C_{1,1}$ 之后发送消息 m；
  \item P2 收到 m，并把接收状态写入 $C_{2,2}$；
  \item P1 故障，回滚到 $C_{1,1}$，该检查点不包含“发送 m”；
  \item 若 P2 仍恢复到 $C_{2,2}$，就出现“接收存在、发送不存在”的 orphan；
  \item 所以恢复必须同时约束多个进程的检查点。
\end{enumerate}
\end{frame}

\sectionframe{二、非协调检查点：自由，但可能产生 Domino Effect}{每个进程独立保存状态，正常运行开销低；恢复时却可能被依赖链拖回很远。}

\begin{frame}{非协调检查点的优点与代价}
\begin{columns}[T]
\column{.48\linewidth}
\blocktitle{优点}
\begin{itemize}
  \item 不需要进程同步；
  \item 正常执行不被阻塞；
  \item 不同进程可按自身负载选择频率。
\end{itemize}
\column{.48\linewidth}
\blocktitle{代价}
\begin{itemize}
  \item 恢复时要搜索一致检查点集合；
  \item 依赖链可能迫使更多进程回滚；
  \item 需要保存多个历史检查点并做垃圾回收。
\end{itemize}
\end{columns}
\vfill
\callout{适用性}{频繁输出提交、恢复速度要求高的系统通常不喜欢完全非协调方案。}
\end{frame}

\begin{frame}{Domino Effect：回滚像骨牌一样传播}
\begin{center}
\resizebox{.97\linewidth}{!}{%
\begin{tikzpicture}[node distance=8mm and 14mm,>=Latex]
  \nodebox[minimum width=2.6cm]{ADSBlue}{p1c}{P1 回滚};
  \nodebox[right=of p1c,minimum width=2.6cm]{ADSOrange}{p2c}{P2 撤销接收};
  \nodebox[right=of p2c,minimum width=2.6cm]{ADSGreen}{p3c}{P3 撤销后续发送};
  \nodebox[right=of p3c,minimum width=2.6cm]{ADSDarkBlue}{p4c}{P4 再次回滚};
  \draw[\arrow,ADSGray] (p1c)--node[above,font=\scriptsize]{消息依赖}(p2c);
  \draw[\arrow,ADSGray] (p2c)--node[above,font=\scriptsize]{因果依赖}(p3c);
  \draw[\arrow,ADSGray] (p3c)--node[above,font=\scriptsize]{继续传播}(p4c);
\end{tikzpicture}
}
\end{center}
\begin{itemize}
  \item 失败只发生在 P1，但恢复影响未失败进程；
  \item 如果依赖链追溯到初始检查点，系统几乎丢失全部进展。
\end{itemize}
\end{frame}

\begin{frame}{依赖跟踪：把“谁依赖谁”记录下来}
\begin{itemize}
  \item 将执行划分为检查点间隔 $I_{i,x}$；
  \item Pj 在间隔 $I_{j,y}$ 收到来自 Pi 的消息时，记录 $I_{i,x}\to I_{j,y}$；
  \item 这些依赖随 $C_{j,y}$ 保存到稳定存储；
  \item 恢复时只需寻找满足依赖闭包的检查点组合，而不是遍历所有历史。
\end{itemize}
\end{frame}

\begin{frame}{恢复线的选择}
\begin{itemize}
  \item 恢复线是每个进程选择的检查点集合；
  \item 它必须没有 orphan，且不能包含发送/接收顺序矛盾；
  \item 越靠后，丢失的计算越少，但越难满足跨进程依赖；
  \item 可以把恢复线看成一个约束求解问题：每个检查点是变量，消息依赖是约束。
\end{itemize}
\end{frame}

\sectionframe{三、协调检查点：共同建立恢复线}{通过协调器和协议让一组局部检查点组成一致全局检查点，从根源上避免 Domino Effect。}

\begin{frame}{阻塞式协调检查点}
\begin{enumerate}
  \item 协调器请求所有进程创建检查点；
  \item 进程保存本地状态后暂停应用消息处理；
  \item 所有进程确认后，检查点集合成为全局恢复线；
  \item 恢复执行，再解除阻塞。
\end{enumerate}
\vfill
\callout{优点/代价}{一致性简单、证明容易；但检查点期间会阻塞计算，延迟敏感服务难以接受。}
\end{frame}

\begin{frame}{非阻塞式协调检查点的难点}
\begin{itemize}
  \item 进程保存检查点后继续运行；
  \item 某进程可能在收到检查点请求前发送消息；
  \item 接收者可能先收到消息，再收到检查点请求；
  \item 如果直接保存接收状态，会出现接收在新检查点、发送在旧检查点的矛盾；
  \item FIFO 信道可以用“检查点标记先行”解决部分问题。
\end{itemize}
\end{frame}

\begin{frame}{FIFO 信道上的标记法}
\begin{center}
\begin{tikzpicture}[x=1cm,y=1cm,>=Latex]
  \draw[ADSBlue,line width=1.2pt,->] (0,1)--(10,1);
  \node[above,font=\scriptsize] at (2,1){checkpoint request};
  \node[above,font=\scriptsize] at (5,1){first post-checkpoint message};
  \draw[ADSOrange,line width=1.2pt,->] (0,0)--(10,0);
  \node[above,font=\scriptsize] at (2,0){P1};\node[above,font=\scriptsize] at (7,0){P2};
  \draw[dashed,ADSGray] (3,-.3)--(3,1.4);
  \node[below,font=\scriptsize] at (3,-.3){P2 先保存检查点};
\end{tikzpicture}
\end{center}
\begin{itemize}
  \item 在 FIFO 信道上，检查点请求先于后续应用消息到达；
  \item 接收者先保存本地状态，再处理第一条 post-checkpoint 消息；
  \item 从而避免接收事件跨过恢复线。
\end{itemize}
\end{frame}

\begin{frame}{FIFO Marker 完整例子：消息落在哪个检查点？}
\small
$P_1$ 记录 $C_1$ 后，在应用消息之前发送 Marker；FIFO 保证同一信道不越过 Marker。
\begin{center}
\begin{tabularx}{.96\linewidth}{>{\bfseries}p{3cm}p{2.8cm}X}
\toprule
情况 & 信道顺序 & $P_2$ 的处理 \\
\midrule
$m$ 在检查点前发送 & $m \rightarrow Marker$ & 先接收 $m$，再记录 $C_2$；$m$ 的发送和接收都属于旧 epoch。 \\
$m'$ 在检查点后发送 & $Marker \rightarrow m'$ & 先记录 $C_2$，再接收 $m'$；$m'$ 的发送和接收都属于新 epoch。 \\
\bottomrule
\end{tabularx}
\end{center}
\begin{enumerate}
  \item $P_1$ 记录本地状态并在所有输出信道发送 Marker；
  \item $P_2$ 第一次收到 Marker 时记录本地状态，再向外转发 Marker；
  \item 每条信道由 Marker 划分为检查点前/后的消息；
  \item 因此不会出现“记录接收，却没有记录对应发送”的 orphan。
\end{enumerate}
\end{frame}

\begin{frame}{Koo--Toueg 两阶段协调思想}
\begin{enumerate}
  \item 发起进程创建 tentative checkpoint，并请求其他进程创建 tentative checkpoint；
  \item 收齐确认后，协调器提交这些 checkpoint 为 permanent；
  \item 若任意进程拒绝或失败，撤销 tentative checkpoint 并重试；
  \item 恢复时从最近一次成功提交的全局检查点开始。
\end{enumerate}
\vfill
\callout{与事务的类比}{tentative/permanent 很像 prepare/commit；一致恢复线本质上是一次分布式提交。}
\end{frame}

\sectionframe{四、通信诱导检查点：把协议元数据放进应用消息}{不额外广播协调消息，而是利用消息上的因果信息决定接收者是否必须强制保存检查点。}

\begin{frame}{自主检查点与强制检查点}
\begin{itemize}
  \item 自主检查点：进程按本地策略自行创建；
  \item 强制检查点：接收消息前，由消息携带的信息触发；
  \item 强制检查点必须先完成，再让应用处理消息内容；
  \item 这样可以推进全局恢复线，同时避免所有进程同步停顿。
\end{itemize}
\end{frame}

\begin{frame}{通信诱导的因果解释}
\begin{itemize}
  \item 消息携带发送者当前 checkpoint index、逻辑时间或依赖摘要；
  \item 接收者比较自己的恢复线与消息要求；
  \item 如果直接处理会让接收状态依赖一个尚未安全保存的发送状态，就先强制 checkpoint；
  \item 该机制把“全局协调”分摊到正常消息流中。
\end{itemize}
\vfill
\callout{代价}{每条消息增加元数据，接收者可能在高峰期频繁强制 checkpoint。}
\end{frame}

\sectionframe{五、日志型恢复：不只保存状态，还保存非确定性}{如果应用满足 Piecewise Deterministic（PWD）假设，保存检查点和事件决定因子即可重放执行。}

\begin{frame}{确定性与非确定性事件}
\begin{itemize}
  \item 进程内部计算通常是确定性的：给定输入，下一步状态唯一；
  \item 收到哪个消息、消息何时到达、随机数和外部 I/O 可能是非确定性的；
  \item 日志记录非确定性事件的 determinant；
  \item 恢复时从检查点重新执行，并按日志重放相同 determinant。
\end{itemize}
\end{frame}

\begin{frame}{No-orphans 条件}
对非确定性事件 $e$，令 $Depend(e)$ 为受其影响的进程集合，$Log(e)$ 为已在内存记录其 determinant 的进程集合，$Stable(e)$ 表示 determinant 已稳定写入。
\[
\neg Stable(e)\Rightarrow Depend(e)\subseteq Log(e)
\]
\begin{itemize}
  \item 如果 determinant 尚未落盘，所有依赖它的进程必须仍能回滚/重放；
  \item 这样不会出现某个进程已观察到事件、但恢复时找不到事件来源的 orphan。
\end{itemize}
\end{frame}

\begin{frame}{悲观日志}
\begin{itemize}
  \item 假设故障可能发生在任意非确定性事件之后；
  \item 事件 determinant 在允许依赖它之前同步写入稳定存储；
  \item 优点：恢复简单、进程之间相互隔离、不会产生 orphan；
  \item 代价：每个非确定性事件都引入稳定存储延迟。
\end{itemize}
\end{frame}

\begin{frame}{乐观日志与因果日志}
\begin{tabularx}{\linewidth}{>{\bfseries}p{2.8cm}XX}
\toprule
乐观日志 & 异步写日志，假设故障很少；恢复时需要追踪因果依赖和多个检查点。 \\
因果日志 & 利用因果依赖保证 no-orphans，平时不强制同步落盘，恢复协议更复杂。 \\
\bottomrule
\end{tabularx}
\vfill
\callout{工程取舍}{悲观方案用运行时延迟换恢复简单；乐观/因果方案用更复杂的恢复和垃圾回收换正常路径性能。}
\end{frame}

\sectionframe{六、AI 训练中的检查点与恢复}{训练 Checkpoint 不是“保存一个文件”，而是保存计算、优化器、数据迭代和通信状态的一致切。}

\begin{frame}{训练 Checkpoint 应保存什么？}
\begin{itemize}
  \item 模型参数、优化器状态、学习率调度器和随机数种子；
  \item 数据集 shard、样本游标和 epoch/step；
  \item 混合精度 scaler、梯度累积和并行组状态；
  \item 版本、代码提交、容器镜像 digest 和硬件拓扑；
  \item 只有这些状态组合一致，恢复后才会沿同一条训练历史继续。
\end{itemize}
\end{frame}

\begin{frame}{分布式训练的全局恢复线}
\begin{center}
\begin{tikzpicture}[node distance=8mm and 14mm,>=Latex]
  \nodebox[minimum width=2.5cm]{ADSBlue}{w1}{Worker 1\\step 100};
  \nodebox[right=of w1,minimum width=2.5cm]{ADSOrange}{w2}{Worker 2\\step 100};
  \nodebox[right=of w2,minimum width=2.5cm]{ADSGreen}{w3}{Worker 3\\step 100};
  \nodebox[below=of w2,minimum width=2.8cm]{ADSDarkBlue}{ck}{Checkpoint\\commit 100};
  \draw[\arrow,ADSGray] (w1)--(ck);\draw[\arrow,ADSGray] (w2)--(ck);\draw[\arrow,ADSGray] (w3)--(ck);
\end{tikzpicture}
\end{center}
\begin{itemize}
  \item 所有 worker 必须对同一个 global step 达成一致；
  \item 只有 Checkpoint manifest 原子提交后，调度器才应把作业标记为可恢复。
\end{itemize}
\end{frame}

\begin{frame}{为什么“每个 worker 各自保存”会出问题？}
\begin{itemize}
  \item Worker 1 保存到 step 100，Worker 2 仍在 step 99；
  \item Worker 1 故障后恢复到 100，Worker 2 恢复到 99；
  \item 梯度、数据游标和优化器状态不再对应同一训练历史；
  \item 需要 barrier、两阶段提交或 manifest 选择一个一致恢复线。
\end{itemize}
\end{frame}

\begin{frame}{Kubernetes 与 Checkpoint 生命周期}
\begin{enumerate}
  \item 训练 Job 创建 worker Pod 和 PVC；
  \item worker 将临时 checkpoint 写入本地或共享存储；
  \item coordinator 验证所有分片、校验和和版本；
  \item 原子写入 manifest，标记 checkpoint 可恢复；
  \item 节点故障后，控制器只从最新完整 manifest 恢复，不读取半成品。
\end{enumerate}
\end{frame}

\begin{frame}{近期案例：PyTorch Distributed Checkpoint / FSDP}
\begin{itemize}
  \item 分布式训练的状态常按 rank 分片保存，而不是由 rank-0 汇总成一个超大文件；
  \item checkpoint 元数据要记录参数、优化器、分片布局、版本和校验信息，恢复时才能检查完整性；
  \item 重新启动的 worker 数量或设备布局可能变化，读取端需要把保存时的分片重新映射到当前 rank；
  \item 因此“文件存在”不等于“恢复线成立”：manifest、分片校验和训练 step 必须共同提交。
\end{itemize}
\vfill
\smallnote{实现名称和 API 会随框架版本变化；本页抽象的是分片状态、元数据和可重分片恢复这三个分布式系统问题。}
\end{frame}

\begin{frame}{课堂互动：判断一个恢复点是否一致}
\begin{enumerate}
  \item P1 的模型参数为 step 200，优化器状态为 199；
  \item P2 的参数和优化器均为 200，但数据游标为 198；
  \item manifest 已写入，但 shard-3 校验和失败；
  \item 讨论：哪些状态可以作为恢复线？需要回滚到哪里？
\end{enumerate}
\vfill
\callout{判断原则}{恢复线必须同时满足版本、数据、通信和持久化完整性，不能只看一个“checkpoint 文件存在”。}
\end{frame}

\sectionframe{七、检查点与回滚完整例子}{用一条带有 orphan、lost 和 delayed 消息的时间线，决定恢复点。}

\begin{frame}{例子：三个进程的消息历史}
\begin{center}
\begin{tikzpicture}[>=Latex,x=0.16cm,y=1cm]
\draw[ADSGray,->] (0,0)--(68,0); \draw[ADSGray,->] (0,-1.2)--(68,-1.2); \draw[ADSGray,->] (0,-2.4)--(68,-2.4);
\node[anchor=east,font=\scriptsize] at (-1,0){P1};\node[anchor=east,font=\scriptsize] at (-1,-1.2){P2};\node[anchor=east,font=\scriptsize] at (-1,-2.4){P3};
\foreach \t in {10,20,30,40,50,60}{\draw[ADSGray] (\t,.08)--(\t,-.08) node[below,font=\tiny]{\t};}
\draw[ADSBlue,very thick] (20,0) circle (2pt) node[above,font=\scriptsize]{send a};
\draw[ADSOrange,very thick] (30,-1.2) circle (2pt) node[above,font=\scriptsize]{recv a};
\draw[ADSGreen,very thick] (40,-1.2) circle (2pt) node[above,font=\scriptsize]{send b};
\draw[ADSGray,very thick] (50,-2.4) circle (2pt) node[below,font=\scriptsize]{recv b};
\draw[\arrow,ADSBlue] (20,-.1)--(30,-1.1);
\draw[\arrow,ADSOrange] (40,-1.3)--(50,-2.3);
\end{tikzpicture}
\end{center}
\begin{itemize}
  \item 若检查点记录了 P2 已接收 a，却没有记录 P1 发送 a，恢复后会出现 orphan；
  \item 若记录 P1 发送 a，却没有记录 P2 接收 a，则 a 可能成为 lost message。
\end{itemize}
\end{frame}

\begin{frame}{选择恢复线：为什么不能各自选最新？}
\begin{enumerate}
  \item P1 的 checkpoint 在事件 25；P2 在事件 35；P3 在事件 45；
  \item 若三者直接各自恢复最新状态，P2 可能依赖 P1 尚未恢复的发送；
  \item 一致恢复线必须是一个一致 cut，不能包含接收却不包含其因果发送；
  \item 需要回退某些较新的局部检查点，或补充日志重放消息。
\end{enumerate}
\end{frame}

\begin{frame}{协调式检查点的两阶段流程}
\begin{enumerate}
  \item coordinator 发送 PREPARE，阻止跨 checkpoint 的新依赖；
  \item 每个进程写本地状态并确认；
  \item coordinator 收齐确认后写 COMMIT manifest；
  \item 进程收到 COMMIT 后把检查点标记为可恢复；
  \item 超时则 ABORT，不能让半完成的检查点冒充最新版本。
\end{enumerate}
\end{frame}

\begin{frame}{非协调检查点：更少阻塞，更多恢复逻辑}
\begin{itemize}
  \item 进程可以独立保存状态，不等待全局 barrier；
  \item 恢复时根据消息日志寻找一致组合；
  \item 需要处理 orphan、lost、duplicate 和重放顺序；
  \item 适合长时间运行、低延迟系统，但恢复算法复杂。
\end{itemize}
\end{frame}

\begin{frame}{日志型恢复：悲观、乐观与因果日志}
\begin{tabularx}{\linewidth}{>{\bfseries}p{3cm}XX}
\toprule
悲观日志 & 消息先持久化再交付。 & 恢复简单，运行时延迟高。 \\
乐观日志 & 先交付，故障后补写日志。 & 运行快，恢复可能回滚更多。 \\
因果日志 & 只记录恢复所需的因果上下文。 & 存储较省，协议和证明更复杂。 \\
\bottomrule
\end{tabularx}
\end{frame}

\begin{frame}{训练作业：半个 checkpoint 为什么危险？}
\begin{itemize}
  \item rank-0 权重已写完，rank-1 权重仍上传中；
  \item manifest 若提前可见，恢复器会读取混合版本；
  \item 正确做法是分片校验、临时目录、原子 rename 和 manifest commit；
  \item 调度器只选择最后一个完整、签名和版本匹配的 manifest。
\end{itemize}
\end{frame}

\begin{frame}{回滚与外部副作用}
\begin{enumerate}
  \item 计算状态可以回滚，已经发送的邮件、支付和模型发布不一定能回滚；
  \item 用 outbox 记录待提交副作用，等 checkpoint/事务达到安全线后再发送；
  \item 每个副作用带幂等键，重复恢复不会重复扣费；
  \item fencing 阻止旧实例继续写入外部系统。
\end{enumerate}
\end{frame}

\begin{frame}{课堂练习：判断四个 checkpoint}
\begin{tabularx}{\linewidth}{>{\bfseries}p{2.3cm}XXXX}
\toprule
检查点 & 参数 & 优化器 & 数据游标 & manifest \\
C1 & 200 & 200 & 200 & 完整 \\
C2 & 200 & 199 & 200 & 完整 \\
C3 & 201 & 201 & 199 & 完整 \\
C4 & 201 & 201 & 201 & shard-3 损坏 \\
\bottomrule
\end{tabularx}
\vfill
\callout{问题}{哪些检查点可直接恢复？哪些必须回退或重算？请说明“完整”不只是文件存在。}
\end{frame}

\begin{frame}{本讲总结}
\begin{enumerate}
  \item 局部检查点只有组合成一致全局检查点，才能安全恢复；
  \item orphan、lost、delayed、duplicate 消息解释了恢复中的主要陷阱；
  \item 非协调、协调、通信诱导和日志型方案分别在运行时开销与恢复复杂度之间取舍；
  \item Checkpoint/Recovery 与分布式提交、因果顺序、幂等和输出提交密切相关；
  \item AI 训练必须把模型、优化器、数据游标、随机性和 manifest 一起纳入恢复线。
\end{enumerate}
\vfill
\callout{下一讲衔接}{恢复依赖“谁还活着、谁可能失联”的判断，下一讲进入 Failure Detector 与 Consensus。}
\end{frame}

\begin{frame}{课后题}
\begin{enumerate}
  \item 画出三个进程的 orphan/lost/delayed 消息例子；
  \item 比较阻塞式协调检查点与通信诱导检查点；
  \item 解释悲观日志为何能避免 no-orphans，乐观日志需要哪些额外恢复信息；
  \item 为 8 卡训练作业设计一个两阶段 checkpoint manifest 协议。
\end{enumerate}
\end{frame}

\end{document}
